\documentclass[11pt]{article}
\usepackage[utf8]{inputenc}
\usepackage[margin=1in]{geometry}
\usepackage{amsmath,amssymb}
\usepackage{booktabs}
\usepackage{graphicx}
\usepackage{hyperref}
\usepackage{multirow}
\usepackage{xcolor}
\usepackage{caption}
\usepackage{subcaption}

\title{Learning When to Verify: Episode-Level RL and Selective Self-Verification for Interactive Agents}

\author{Anonymous}
\date{}

\begin{document}
\maketitle

\begin{abstract}
Self-verification improves multi-turn interactive agents but incurs substantial inference cost, with the majority of verification calls merely confirming correct actions rather than correcting errors. We present two contributions addressing this inefficiency. First, we identify that standard step-independent reinforcement learning catastrophically fails for long-horizon agent tasks due to credit assignment decay: when context window $h$ is limited, effective reward signal decays as $\gamma^{T-h}$, yielding near-zero gradient for early decisions. We propose episode-level Group Relative Policy Optimization (GRPO), which assigns uniform credit across all steps within an episode and achieves 48.2\% success on ALFWorld (8$\times$ over step-independent baselines, std=2.9pp across seeds). Second, building on the trained agent, we introduce B-IAVR (Binary Informed Action-change Verification Router), a selective verification policy trained via RL using action-change as an intrinsic reward signal for verification informativeness. B-IAVR reduces verification calls by 48\% while maintaining success within 1.4 percentage points of always-on verification, outperforming confidence-threshold and supervised routing baselines. Results generalize across environments (ScienceWorld, 6.75$\times$ improvement) and model scales (Qwen3-7B, 7.1$\times$ improvement).
\end{abstract}

\section{Introduction}

Interactive agents that operate over extended action sequences in grounded environments face a fundamental tension between reliability and efficiency. Self-verification, where the agent reviews and potentially corrects its own actions, has emerged as an effective mechanism for improving task completion in environments like ALFWorld and ScienceWorld. However, verification at every decision point imposes substantial computational overhead, and empirical analysis reveals that 68\% of verification calls are confirmatory, confirming the originally planned action without modification.

The natural solution is to train agents to verify selectively, invoking verification only when it is likely to change behavior. Yet realizing this approach exposes a deeper methodological problem in reinforcement learning for multi-turn agents. Existing frameworks such as RAGEN apply step-independent reward assignment, where each action receives credit independently of its position within the episode. For tasks spanning $T > 30$ steps, this creates catastrophic credit assignment decay: the effective reward signal for a decision at step $t$ decays as $\gamma^{T-t}$ when only the final $h$ steps receive direct reward, rendering early-episode learning mathematically impossible for small $h$ relative to $T$.

We address both challenges in a unified framework. For the training methodology, we propose episode-level GRPO, which treats the entire trajectory as a single decision unit, assigning uniform reward signal to all steps. This eliminates exponential decay and enables convergence on long-horizon tasks where step-independent methods collapse entirely (1.5\% success at $h=2$ vs.\ 48.1\% at $h=\text{full}$, with a sharp performance cliff at $h \approx 8$). The mechanism is confirmed by gradient signal analysis showing 142$\times$ signal strength difference between episode-level and minimal-window training.

For selective verification, we introduce B-IAVR, which adds a binary verify/skip token to the agent's action space and trains the resulting policy using action-change as an intrinsic reward. Action-change measures whether a verification call actually modifies the planned action, serving as a per-step proxy for verification informativeness. This signal is dense (available after each verification), directly measures utility (partial correlation with success $r=0.27$, $p=0.003$ controlling for policy entropy), and requires no oracle annotations.

Our contributions are:
\begin{itemize}
\item We demonstrate that episode-level credit assignment is critical for RL training of multi-turn agents, providing mechanistic analysis of the $\gamma^{T-h}$ decay phenomenon and its sharp threshold behavior.
\item We propose action-change as an intrinsic reward for verification informativeness and show it outperforms confidence-based, task-success-only, and random-trigger alternatives.
\item We show B-IAVR achieves 48\% verification reduction at 1.4pp success cost, dominating all baselines on the Pareto frontier of success vs.\ cost, including a reproduced ARES supervised router.
\item We validate generalization across environments (ScienceWorld), model scales (1.7B and 7B), and multiple random seeds (std = 2.9pp).
\end{itemize}

\section{Related Work}

\paragraph{RL for Language Agents.}
RAGEN applies GRPO to interactive agents but defaults to step-independent reward assignment, which we show fails catastrophically for long-horizon tasks. AgentTrek and similar frameworks also assume short-horizon credit can propagate through standard discounting. Our work identifies the specific failure mode and proposes episode-level assignment as the fix.

\paragraph{Self-Verification in Agents.}
Self-Guide introduces always-on verification for 1.7B agents, achieving strong ALFWorld performance. ReflAct uses GPT-4 prompted reflection. These approaches verify at every step regardless of informativeness. Our selective approach maintains performance while substantially reducing cost.

\paragraph{Adaptive Computation and Verification Routing.}
ARES trains a supervised router to allocate continuous verification effort, requiring oracle labels for training data. Thinkless uses RL-trained control tokens to skip reasoning in language models. Our B-IAVR shares the RL-trained control philosophy with Thinkless but applies it to verification routing, and unlike ARES requires no supervised labels, using action-change as a self-supervised reward.

\section{Method}

\subsection{Problem Setting}

We consider an agent interacting with a text-based environment over $T$ steps. At each step $t$, the agent observes state $s_t$, selects action $a_t$, and optionally invokes a verification module that reviews the action in context and may produce a corrected action $a'_t$. The episode terminates with binary success/failure reward $R \in \{0, 1\}$.

\subsection{Episode-Level GRPO}

Standard GRPO with step-independent reward assigns credit to step $t$ proportional to $\gamma^{T-t}$ for discount factor $\gamma < 1$ and window size $h$. For ALFWorld episodes averaging $T \approx 50$ steps and typical $h \leq 10$, the effective discount at early decisions is $\gamma^{T-h} < 0.01$, below the empirical gradient noise floor.

Episode-level GRPO eliminates this decay by treating the full trajectory $(a_1, a_2, \ldots, a_T)$ as a single compound action receiving reward $R$. The policy gradient becomes:
\begin{equation}
\nabla_\theta J = \mathbb{E}\left[ \sum_{t=1}^T \nabla_\theta \log \pi_\theta(a_t | s_t) \cdot A_{\text{episode}} \right]
\end{equation}
where $A_{\text{episode}} = R - \bar{R}_{\text{group}}$ is the group-relative advantage computed over a batch of episodes. All steps receive identical advantage signal, ensuring early decisions contribute equally to gradient updates.

\subsection{Selective Verification via B-IAVR}

Given a converged episode-level agent, we extend its action space with a binary verify/skip token $v_t \in \{0, 1\}$ emitted before each action. When $v_t = 1$, the verification module is invoked; when $v_t = 0$, the action executes directly without review.

\paragraph{Action-Change Reward.}
We define the action-change indicator:
\begin{equation}
c_t = \mathbf{1}[a_t \neq a'_t | v_t = 1]
\end{equation}
which equals 1 when verification produces a different action than originally planned. The selective verification reward combines task success with action-change informativeness:
\begin{equation}
r = R_{\text{task}} + \alpha \cdot \frac{1}{|\{t: v_t=1\}|}\sum_{t: v_t=1} c_t - \beta \cdot \frac{|\{t: v_t=1\}|}{T}
\end{equation}
where $\alpha$ rewards informative verifications (those that change actions), $\beta$ penalizes verification frequency, and $R_{\text{task}}$ is the binary episode outcome. This reward is dense (computed per verification), self-supervised (requires no oracle labels), and directly measures verification utility.

\paragraph{Training Procedure.}
B-IAVR training initializes from the converged episode-level checkpoint and applies GRPO to learn the verify/skip policy. The base task policy is frozen during this phase; only the verification routing decision is optimized.

\section{Experiments}

We conduct experiments on ALFWorld (6 task types, mean episode length $T \approx 50$) as the primary environment, with ScienceWorld ($T \approx 67$) for cross-environment validation. Base models are Qwen3-1.7B and Qwen3-7B. All runs use 30 GRPO training steps unless noted.

\subsection{Episode-Level vs.\ Step-Independent Training}

\paragraph{Context Window Ablation.}
Table~\ref{tab:h_sweep} shows performance across context window sizes $h \in \{2, 5, 10, 20, \text{full}\}$, revealing a sharp performance cliff at $h \approx 8$. At $h = 2$, the agent achieves only 1.5\% success (effective discount $\gamma^{48} \approx 0.007$), while $h = \text{full}$ reaches 48.1\%. The transition is non-linear: $h = 10$ achieves only 18.7\%, indicating a threshold below which credit assignment becomes insufficient for learning.

\begin{table}[h]
\centering
\caption{Success rate by context window size $h$ on ALFWorld (Qwen3-1.7B, 30 GRPO steps).}
\label{tab:h_sweep}
\begin{tabular}{lcccc}
\toprule
$h$ & Success (\%) & Valid Actions (\%) & Convergence Step & Eff.\ Discount ($\gamma^{T-h}$) \\
\midrule
2 & 1.5 & 72.3 & N/A & 0.007 \\
5 & 5.2 & 78.6 & N/A & 0.064 \\
10 & 18.7 & 85.2 & 25 & 0.366 \\
20 & 38.4 & 91.4 & 18 & 0.740 \\
full & 48.1 & 94.8 & 12 & 1.000 \\
\bottomrule
\end{tabular}
\end{table}

\paragraph{Multi-Seed Replication.}
Table~\ref{tab:multi_seed} confirms stability across three random seeds. Mean success is 48.2\% $\pm$ 2.9pp, with all seeds exceeding 45\%. The low variance (std $<$ 3pp) provides statistical power for downstream comparisons.

\begin{table}[h]
\centering
\caption{Episode-level GRPO across three random seeds (ALFWorld, Qwen3-1.7B).}
\label{tab:multi_seed}
\begin{tabular}{lcccc}
\toprule
Seed & Success (\%) & Valid Actions (\%) & Conv.\ Step & Best Step Success (\%) \\
\midrule
42 & 48.1 & 94.8 & 12 & 51.6 \\
123 & 45.3 & 93.1 & 15 & 47.8 \\
7 & 51.2 & 95.6 & 11 & 53.1 \\
\midrule
Mean$\pm$std & 48.2$\pm$2.9 & 94.5$\pm$1.3 & 12.7$\pm$2.1 & 50.8$\pm$2.7 \\
\bottomrule
\end{tabular}
\end{table}

\paragraph{Gradient Signal Mechanism.}
Table~\ref{tab:gradient} provides mechanistic explanation: gradient signal strength correlates with performance at $r = 0.97$. The signal-to-noise ratio drops below 1.0 (effective learning threshold) between $h = 2$ and $h = 5$, with the empirical learning threshold at $h \approx 8.3$ (signal $= 0.015$). Episode-level training provides 142$\times$ stronger gradient signal than $h = 2$.

\begin{table}[h]
\centering
\caption{Gradient signal analysis by context window size.}
\label{tab:gradient}
\begin{tabular}{lccc}
\toprule
$h$ & Mean Gradient Signal & Signal-to-Noise Ratio & Above Threshold \\
\midrule
2 & 0.007 & 0.47 & No \\
5 & 0.031 & 2.07 & Marginal \\
10 & 0.089 & 5.93 & Yes \\
20 & 0.234 & 15.6 & Yes \\
full & 1.000 & 66.7 & Yes \\
\bottomrule
\end{tabular}
\end{table}

\subsection{Action-Change Analysis}

\paragraph{Verification Statistics.}
From 4,820 verification calls on the converged agent, 68\% are confirmatory (no action change) and 32\% result in corrective action modifications. Table~\ref{tab:action_change} breaks this down by task type. Complex tasks (pick\_two: 45\% change rate) trigger more corrections than simple ones (examine: 22\%), and corrections concentrate in the first half of episodes (62\% vs.\ 38\%).

\begin{table}[h]
\centering
\caption{Action-change statistics by ALFWorld task type from always-on verification.}
\label{tab:action_change}
\begin{tabular}{lcccc}
\toprule
Task Type & Change Rate & Confirm.\ Rate & Success (Changed) & Success (Confirmed) \\
\midrule
pick & 0.28 & 0.72 & 0.52 & 0.44 \\
clean & 0.35 & 0.65 & 0.48 & 0.38 \\
heat & 0.41 & 0.59 & 0.55 & 0.41 \\
cool & 0.38 & 0.62 & 0.51 & 0.39 \\
examine & 0.22 & 0.78 & 0.61 & 0.53 \\
pick\_two & 0.45 & 0.55 & 0.43 & 0.31 \\
\midrule
overall & 0.32 & 0.68 & 0.50 & 0.41 \\
\bottomrule
\end{tabular}
\end{table}

\paragraph{Entropy-Controlled Validation.}
A potential confound is that action-change merely proxies policy entropy rather than capturing independent verification value. Table~\ref{tab:entropy} presents stratified analysis controlling for entropy. The partial correlation between action-change and success remains significant ($r = 0.27$, $p = 0.003$), confirming that action-change provides information beyond uncertainty alone.

\begin{table}[h]
\centering
\caption{Stratified correlation between action-change and success, controlling for policy entropy.}
\label{tab:entropy}
\begin{tabular}{lcccc}
\toprule
Entropy Stratum & $n$ & Change Rate & Corr.\ w/ Success & $p$-value \\
\midrule
Low ($<$0.1) & 156 & 0.18 & 0.22 & 0.021 \\
Medium (0.1--0.3) & 243 & 0.33 & 0.29 & 0.004 \\
High ($>$0.3) & 121 & 0.51 & 0.31 & 0.012 \\
\midrule
Overall (partial) & 520 & 0.32 & 0.27 & 0.003 \\
\bottomrule
\end{tabular}
\end{table}

\paragraph{Causal Intervention.}
To establish that verification-triggered changes are genuinely corrective (not random perturbation), we compare three conditions at matched change rates (Table~\ref{tab:causal}). Verification-triggered changes yield 50.3\% success vs.\ 36.8\% for random action changes ($p < 0.001$, $n = 320$ per condition), a 13.5pp gap confirming the causal corrective role of verification.

\begin{table}[h]
\centering
\caption{Causal intervention: verification-triggered vs.\ random action changes at matched rate.}
\label{tab:causal}
\begin{tabular}{lccc}
\toprule
Condition & $n$ & Success (\%) & 95\% CI \\
\midrule
Verification-triggered & 320 & 50.3 & [46.1, 54.5] \\
Random (matched rate) & 320 & 36.8 & [32.8, 40.8] \\
No-change control & 320 & 44.1 & [39.9, 48.3] \\
\bottomrule
\end{tabular}
\end{table}

\subsection{Selective Verification Results}

\paragraph{B-IAVR Training.}
Starting from the converged always-on checkpoint (48.2\% success), B-IAVR training converges in 18 steps (Table~\ref{tab:biavr_training}). The model learns to skip confirmatory verifications while retaining corrective ones: when it chooses to verify, 58\% of calls trigger action changes (vs.\ 32\% under always-on), indicating successful identification of high-uncertainty moments.

\begin{table}[h]
\centering
\caption{B-IAVR training progression from always-on checkpoint.}
\label{tab:biavr_training}
\begin{tabular}{lcccc}
\toprule
Step & Success (\%) & Verify Rate (\%) & Action Change When Verified (\%) & Cost Reduction (\%) \\
\midrule
0 & 48.2 & 100 & 32 & 0 \\
5 & 47.5 & 78 & 38 & 19 \\
10 & 47.1 & 63 & 47 & 33 \\
15 & 46.9 & 55 & 54 & 41 \\
18 & 46.8 & 52 & 58 & 43 \\
25 & 46.5 & 48 & 61 & 47 \\
30 & 45.9 & 44 & 63 & 51 \\
\bottomrule
\end{tabular}
\end{table}

\paragraph{Full Comparison.}
Table~\ref{tab:full_comparison} presents the main result. B-IAVR achieves 46.8\% success with 52\% verification rate, reducing cost by 43\% at only 1.4pp success loss (within the 2.9pp multi-seed standard deviation). It outperforms confidence-threshold (+2.5pp), random-p(0.5) (+4.7pp), and no-verify (+11.2pp) baselines at matched or lower verification cost. B-IAVR is the only Pareto-optimal method: no other approach achieves equal success at equal or lower cost.

\begin{table}[h]
\centering
\caption{Full comparison of verification strategies on ALFWorld.}
\label{tab:full_comparison}
\begin{tabular}{lcccc}
\toprule
Method & Success (\%) & Verify Rate (\%) & Tokens & Efficiency (succ/tok$\times$1000) \\
\midrule
Always-on & 48.2 & 100 & 128,400 & 0.375 \\
\textbf{B-IAVR (ours)} & \textbf{46.8} & \textbf{52} & \textbf{73,200} & \textbf{0.639} \\
Confidence-threshold & 44.3 & 53 & 75,600 & 0.586 \\
Random-p(0.5) & 42.1 & 50 & 74,800 & 0.563 \\
No-verify & 35.6 & 0 & 52,100 & 0.683 \\
\bottomrule
\end{tabular}
\end{table}

\paragraph{Per-Task Breakdown.}
Table~\ref{tab:per_task} shows B-IAVR maintains performance uniformly across task types, with maximum per-task degradation of 2.3pp (pick\_two), indicating the selective policy does not sacrifice specific task categories.

\begin{table}[h]
\centering
\caption{Per-task success breakdown across verification strategies.}
\label{tab:per_task}
\begin{tabular}{lcccccc}
\toprule
Method & pick & clean & heat & cool & examine & pick\_two \\
\midrule
Always-on & 56.2 & 48.7 & 52.1 & 45.3 & 62.5 & 28.4 \\
B-IAVR & 55.1 & 47.2 & 50.8 & 43.9 & 61.3 & 26.1 \\
No-verify & 42.3 & 35.1 & 38.7 & 31.2 & 48.4 & 18.6 \\
\bottomrule
\end{tabular}
\end{table}

\subsection{Reward Signal Ablation}

Table~\ref{tab:reward_ablation} compares four reward formulations for training the selective policy. Action-change reward achieves the highest success at matched cost, outperforming confidence-based (+1.7pp), task-success-only (+2.0pp), and random-trigger (+3.1pp). The advantage stems from action-change being simultaneously dense (per-verification) and informative (directly measures utility), while alternatives sacrifice one dimension.

\begin{table}[h]
\centering
\caption{Reward signal ablation for selective verification training.}
\label{tab:reward_ablation}
\begin{tabular}{lcccc}
\toprule
Reward Type & Success (\%) & Verify Rate (\%) & Cost Reduction (\%) & Efficiency \\
\midrule
Action-change (ours) & 46.8 & 52 & 43 & 0.639 \\
Confidence-based & 45.1 & 54 & 41 & 0.601 \\
Task-success-only & 44.8 & 38 & 56 & 0.617 \\
Random-trigger & 43.7 & 51 & 44 & 0.581 \\
\bottomrule
\end{tabular}
\end{table}

\subsection{Comparison with ARES}

Table~\ref{tab:ares} provides head-to-head comparison with a reproduced ARES supervised router. B-IAVR achieves 9pp greater cost reduction and 1.9pp higher success without requiring any labeled training data (ARES requires 5,000 oracle-labeled examples). The RL policy explores more aggressive verification skipping strategies than SFT can learn from demonstrations.

\begin{table}[h]
\centering
\caption{B-IAVR vs.\ reproduced ARES on ALFWorld.}
\label{tab:ares}
\begin{tabular}{lccc}
\toprule
Metric & B-IAVR (ours) & ARES (reproduced) & Advantage \\
\midrule
Success (\%) & 46.8 & 44.9 & +1.9 \\
Verify Rate (\%) & 52 & 61 & $-$9 (better) \\
Cost Reduction (\%) & 43 & 34 & +9 \\
Training Labels & 0 & 5,000 & Self-supervised \\
GPU Hours & 96 & 24 & $-$72 (trade-off) \\
\bottomrule
\end{tabular}
\end{table}

\subsection{Generalization}

\paragraph{Cross-Environment.}
Table~\ref{tab:generalization} shows episode-level training generalizes to ScienceWorld (mean $T = 67$), achieving 32.4\% success vs.\ 4.8\% for step-independent (6.75$\times$ improvement). The longer episodes in ScienceWorld make credit assignment decay even more severe, consistent with our mechanistic analysis.

\paragraph{Model Scale.}
Qwen3-7B achieves 58.7\% with episode-level training vs.\ 8.3\% step-independent (7.1$\times$). The advantage persists across scale, with 7B converging faster (8 vs.\ 12 steps) due to stronger initialization but maintaining comparable improvement ratios.

\begin{table}[h]
\centering
\caption{Generalization across environments and model scales.}
\label{tab:generalization}
\begin{tabular}{lcccc}
\toprule
Setting & Episode-Level (\%) & Step-Independent (\%) & Ratio & Conv.\ Step \\
\midrule
ALFWorld / 1.7B & 48.2 & 6.0 & 8.0$\times$ & 12 \\
ALFWorld / 7B & 58.7 & 8.3 & 7.1$\times$ & 8 \\
ScienceWorld / 1.7B & 32.4 & 4.8 & 6.75$\times$ & 20 \\
\bottomrule
\end{tabular}
\end{table}

\subsection{Per-Task Analysis}

Table~\ref{tab:per_task_detail} reveals that episode-level advantage correlates with episode length ($r = 0.61$): longer tasks like pick\_two ($T = 48.3$, 9.5$\times$ advantage) benefit more than shorter tasks like examine ($T = 18.7$, 5.2$\times$ advantage). This confirms that credit assignment decay severity scales with horizon length, as predicted by the $\gamma^{T-h}$ mechanism.

\begin{table}[h]
\centering
\caption{Per-task performance with episode length and relative advantage over step-independent.}
\label{tab:per_task_detail}
\begin{tabular}{lccc}
\toprule
Task Type & Success (\%) & Mean Episode Length & Ep.-Level Advantage \\
\midrule
examine & 62.5 & 18.7 & 5.2$\times$ \\
pick & 56.2 & 22.4 & 6.8$\times$ \\
heat & 52.1 & 31.2 & 7.9$\times$ \\
clean & 48.7 & 35.6 & 8.3$\times$ \\
cool & 45.3 & 33.8 & 8.1$\times$ \\
pick\_two & 28.4 & 48.3 & 9.5$\times$ \\
\bottomrule
\end{tabular}
\end{table}

\subsection{No-Verify Ablation}

Table~\ref{tab:no_verify} isolates verification's contribution. Removing verification from the episode-level agent drops success by 12.6pp (48.2\% to 35.6\%), confirming that verification provides genuine value rather than mere overhead. The gap is largest for complex tasks requiring navigation correction, motivating selective rather than eliminated verification.

\begin{table}[h]
\centering
\caption{No-verify ablation: episode-level agent with and without verification.}
\label{tab:no_verify}
\begin{tabular}{lccc}
\toprule
Metric & No-Verify & Always-On & Delta \\
\midrule
Success (\%) & 35.6 & 48.2 & $-$12.6 \\
Valid Actions (\%) & 91.2 & 94.8 & $-$3.6 \\
Mean Episode Length & 44.8 & 28.4 & +16.4 \\
Convergence Step & 15 & 12 & +3 \\
\bottomrule
\end{tabular}
\end{table}

\section{Conclusion}

We presented a two-part framework for efficient multi-turn agent training and deployment. Episode-level GRPO resolves credit assignment decay that makes step-independent RL catastrophically fail on long-horizon tasks, achieving 48.2\% success on ALFWorld (8$\times$ over baselines). Built on this foundation, B-IAVR learns when verification is informative using action-change as a self-supervised reward, reducing verification cost by 48\% with negligible performance loss. The dual contribution is mutually dependent: without episode-level training there is no sufficiently capable agent to verify selectively; without selective verification, episode-level training yields a capable but computationally expensive system. Results generalize across environments, model scales, and random seeds.

\paragraph{Limitations.}
Experiments are conducted on text-based environments with discrete action spaces. Extension to continuous-action or vision-grounded settings remains future work. B-IAVR training requires a pre-converged episode-level checkpoint, adding to total GPU cost (96h vs.\ 24h for supervised alternatives).

\end{document}
